\documentclass[11pt]{amsart}
\usepackage{amsmath,amssymb,amsthm,url}
\newtheorem{theorem}{Theorem}\newtheorem{lemma}[theorem]{Lemma}\newtheorem{proposition}[theorem]{Proposition}
\theoremstyle{remark}\newtheorem{remark}[theorem]{Remark}
\title{Conserved charges in linearly reinforced graph urns}
\author{Yitzchak Shmalo}
\begin{document}
\begin{abstract}
In the graph-based P\'olya urn with linear reinforcement (WARM with $\alpha=1$), every vector $\sigma$ in the kernel of the
unsigned vertex--edge incidence matrix yields an exact martingale $\sum_e\sigma_e H(N_e-1)$, where $H$ is the harmonic
number. These charges are the flat directions of the Lyapunov function, and they determine which equilibrium the process
selects. On the even cycle $C_{L}$ the limit is $(1+u,1-u,\dots)$ with $L\,\mathrm{artanh}(u)$ equal to the terminal value
of the charge, and $\sqrt{L}\,\mathrm{artanh}(u)$ converges in law to a centred Gaussian as $L\to\infty$. For large even
cycles this describes the law of the random limit, in the continuous-time form of the question raised by Chen and Lucas
and by Lima for the discrete-time urn.
\end{abstract}
\maketitle

\section{The model and the result}
In the urn with graph-based competition of van der Hofstad, Holmes, Kuznetsov and Ruszel \cite{HHKR} (WARM), each edge
of a graph carries a count, each vertex carries a Poisson clock, and when a vertex rings it reinforces one of its edges
with probability proportional to $N_e^\alpha$. For $\alpha>1$ the counts localise; at $\alpha=1$ the process sits at the
boundary of that regime. The discrete-time urn, in which every edge adds a ball to one of its two endpoints at each
step, was introduced by Bena\"im, Benjamini, Chen and Lima \cite{BBCL}. At $\alpha=1$ it was studied by Chen and Lucas \cite{CL}, who proved convergence on regular bipartite graphs to a
random point of a segment of equilibria and left the law of that point open, and by Lima \cite{Lima}, who completed the
convergence proof for balanced bipartite graphs and asked for the distribution of the limit (his Question~1). On the
infinite line, Kleptsyn conjectures that the normalised counts of the critical process converge to $1$
\cite[Section~5]{AletOP}.

The random limit is fixed by conserved charges. For every $\sigma$ in the kernel of the unsigned
incidence matrix, $\sum_e\sigma_eH(N_e-1)$ is an exact martingale (Theorem~\ref{thm:mart}); up to a vanishing error it
equals $\sum_e\sigma_e\log(N_e/t)$, which is conserved by the mean-field flow. On the even cycle $C_L$ the kernel is
one-dimensional and the limit is $(1+u_\infty,1-u_\infty,\dots)$ with $L\,\mathrm{artanh}(u_\infty)$ equal to the terminal
value $M_\infty$ of the charge (Theorem~\ref{thm:conv}). The quadratic variation of $M$ is a sum of local terms, and
a law of large numbers for it gives $\sqrt L\,\mathrm{artanh}(u_\infty)\Rightarrow\mathcal N(0,c)$, where $c$ is the
expected quadratic variation per vertex of the process on $\mathbb{Z}$ (Theorem~\ref{thm:clt}). Section~\ref{sec:num}
reports simulations: the predicted variance on cycles, random limits on $C_4$ and $K_4$ whose dimension matches the
kernel, and the slow decay of the alternating envelope on long cycles that the charges predict for the line.

\section{The charge martingales}
Let $G=(V,E)$ be a finite graph. Each vertex carries an independent rate-one Poisson clock; when the clock of $v$ rings,
one edge $e\ni v$ is chosen with probability $N_e/S_v$, where $S_v=\sum_{f\ni v}N_f$, and $N_e$ increases by one. All
counts start at $1$. Write $X_e(t)=N_e(t)/t$ and let $B$ be the unsigned incidence matrix, $B_{ve}=1$ when $v\in e$.

\begin{theorem}\label{thm:mart}
Let $\sigma\in\mathbb{R}^E$ satisfy $\sum_{e\ni v}\sigma_e=0$ for every $v$. Then
\[ M_\sigma(t)=\sum_{e\in E}\sigma_e\,H(N_e(t)-1),\qquad H(n)=\sum_{k=1}^n\frac1k, \]
is a martingale with $M_\sigma(0)=0$, jumps bounded by $\max_e|\sigma_e|$, and predictable quadratic variation
\[ \langle M_\sigma\rangle_t=\int_0^t\sum_{v\in V}\frac{1}{S_v(s)}\sum_{e\ni v}\frac{\sigma_e^2}{N_e(s)}\,ds . \]
\end{theorem}
\begin{proof}
When $N_e$ moves from $n$ to $n+1$, $H(N_e-1)$ increases by $1/n$. A ring at $v$ therefore changes $M_\sigma$ by
$\sigma_e/N_e$ with probability $N_e/S_v$, for each $e\ni v$. The conditional mean of the change is
$S_v^{-1}\sum_{e\ni v}\sigma_e=0$, and its conditional second moment is $S_v^{-1}\sum_{e\ni v}\sigma_e^2/N_e$. The jump
rates are bounded by $|V|$ and the jumps by $\max_e|\sigma_e|$, so the compensated sums are martingales.
\end{proof}

The space of such $\sigma$ is $\ker B$, of dimension $|E|-|V|+\beta$, with $\beta$ the number of bipartite components.
It is nonzero exactly when some component contains an even cycle or two odd cycles. Summing the vertex constraints
gives $\sum_e\sigma_e=0$ for $\sigma\in\ker B$, and $H(n-1)=\psi(n)+\gamma$, so
\begin{equation}\label{eq:link}
\sum_e\sigma_e\log X_e(t)=M_\sigma(t)+\sum_e\sigma_e\bigl(\log N_e(t)-\psi(N_e(t))\bigr),
\end{equation}
with $\log n-\psi(n)=\tfrac1{2n}+O(n^{-2})$. The charges $\sum_e\sigma_e\log X_e$ are therefore martingales up to a term
that vanishes as every count grows. In the mean-field equation $\dot x_e=x_e\sum_{v\in e}S_v^{-1}-x_e$ (time $\log t$)
they are exactly conserved: $\frac{d}{d\tau}\sum_e\sigma_e\log x_e=\sum_v S_v^{-1}\sum_{e\ni v}\sigma_e-\sum_e\sigma_e=0$.

On the even cycle $C_L$ (edges indexed mod $L$, vertex $v$ between edges $v-1$ and $v$) the kernel is spanned by
$\sigma_e=(-1)^e$, and since $S_v=N_{v-1}+N_v$,
\[ \frac{1}{S_v}\Bigl(\frac1{N_{v-1}}+\frac1{N_v}\Bigr)=\frac{1}{N_{v-1}N_v},\qquad
   \langle M\rangle_t=\int_0^t\sum_{v}\frac{ds}{N_{v-1}(s)N_v(s)} .\]

\section{Growth of every edge}
\begin{lemma}\label{lem:growth}
Let $d$ be the maximum degree and $\lambda<2/(d+1)$. Almost surely $\liminf_{t\to\infty}N_e(t)/t^{\lambda}>0$ for every
edge $e$. In particular every count tends to infinity and $\langle M_\sigma\rangle_\infty<\infty$.
\end{lemma}
\begin{proof}
Let $F_w(t)$ be the number of rings of $w$ by time $t$. A ring of $w$ adds to $S_v$ only if $w=v$ or $w\sim v$, and then
by one, so $S_v(t)\le d+\sum_{w\in\{v\}\cup N(v)}F_w(t)$. Since $F_w(t)/t\to1$ a.s., for $\varepsilon>0$ there is an
a.s. finite $T$ with $S_v(t)\le(d+1)(1+\varepsilon)t$ for all $v$ and $t\ge T$. For $e=(u,v)$ the jump rate of $N_e$ is
$N_e(S_u^{-1}+S_v^{-1})\ge\lambda' N_e/t$ on $[T,\infty)$, with $\lambda'=2/((d+1)(1+\varepsilon))$. In the time
$\tau=\log t$ this is a pure birth process with rate at least $\lambda'n$ at height $n$, so by thinning it dominates a
Yule process of rate $\lambda'$ started at time $\log T$ from $N_e(T)\ge1$, and $e^{-\lambda'\tau}Y(\tau)$ has a positive
a.s. limit. Taking $\varepsilon$ small gives the first claim. For the second, every ring of $v$ counts in $S_v$, so
$S_v(t)\ge F_v(t)$, and the rate in Theorem~\ref{thm:mart} is at most
$d\,\|\sigma\|_\infty^2/(F_v(t)\min_{e\ni v}N_e(t))=O(t^{-1-\lambda})$.
\end{proof}

\section{Convergence on the even cycle}
Let $L\ge4$ be even. In the time $\tau=\log t$ the normalised counts follow, up to a martingale error of order
$t^{-1/2}$, the field
\begin{equation}\label{eq:ode}
\frac{dx_e}{d\tau}=x_e\,\partial_e\mathcal{L}(x),\qquad
\mathcal{L}(x)=\sum_{v}\log S_v(x)-\sum_e x_e,
\end{equation}
where $\partial_e\mathcal{L}=S_u^{-1}+S_v^{-1}-1$ for $e=(u,v)$, and $\frac{d}{d\tau}\mathcal{L}=\sum_e x_e(\partial_e\mathcal{L})^2\ge0$ \cite{HHKR}. Also
$\frac{d}{d\tau}\sum_e x_e=L-\sum_e x_e$.

\begin{lemma}\label{lem:eq}
The equilibria of \eqref{eq:ode} in $\{\sum_e x_e=L,\ S_v>0\ \forall v\}$ are
\begin{enumerate}
\item the segment $\Sigma=\{x^{(u)}:u\in[-1,1]\}$, $x^{(u)}_e=1+(-1)^eu$, on which $S_v\equiv2$;
\item finitely many points, each with an edge $e$ such that $x_e=0$ and $\partial_e\mathcal{L}(x)\ge 1/(2L)$.
\end{enumerate}
\end{lemma}
\begin{proof}
Let $x$ be an equilibrium. A vertex with both edges at zero has $S_v=0$, which is excluded. If every edge is positive,
$1/S_u+1/S_v=1$ along every edge of the cycle, so $S$ takes two values $s,s'$ alternately; summing,
$\sum_v S_v=2\sum_e x_e=2L$ gives $s+s'=4$, hence $s=s'=2$ and $x\in\Sigma$. Otherwise the positive edges form paths
$w_0w_1\cdots w_k$ separated by single zero edges, and along such a path $S(w_i)$ alternates between $s$ and
$s'=s/(s-1)$. Since the edges outside the path at $w_0$ and $w_k$ are zero, $x_1=S(w_0)=s$, and solving along the path
gives $x_{2i}=id$, $x_{2i+1}=s-id$ with $d=s'-s$. If $k$ is odd the end condition $x_k=S(w_k)=s'$ forces $d=0$, hence
$k=1$ and $x_1=2$. If $k$ is even it forces $d=2s/k$, and $1/s+1/s'=1$ gives $s=(2k+2)/(k+2)$; all $x_j$ are then
positive. So every block is either a single edge of weight $2$ with $S=2$ at both ends, or an even block of length $k$
with $S=(2k+2)/(k+2)$ at both ends, and the configuration is determined by the block lengths. If all blocks have length
one, $x$ alternates $2,0$ and lies in $\Sigma$ ($u=\pm1$). Otherwise some even block of length $k\le L-1$ meets a zero edge
$e$ whose other end has $S\in\{2,(2k'+2)/(k'+2)\}$, so
$\partial_e\mathcal{L}\ge \frac{k+2}{2k+2}+\frac12-1=\frac{1}{2k+2}\ge\frac1{2L}$, the smallest value occurring when
the other end has $S=2$.
\end{proof}

\begin{theorem}\label{thm:conv}
On $C_L$, $L$ even, almost surely $X(t)\to x^{(u_\infty)}$ with $|u_\infty|<1$ and
\[ L\,\mathrm{artanh}(u_\infty)=M_\infty:=\lim_{t\to\infty}\sum_e(-1)^eH(N_e(t)-1). \]
\end{theorem}
\begin{proof}
By Lemma~\ref{lem:growth}, $S_v(t)\ge F_v(t)$, so $X(t)$ eventually stays in the compact set
$K=\{\sum x_e\le 2L,\ S_v\ge1/2\}$, on which \eqref{eq:ode} is smooth. The jumps of $X$ are of size $1/t$ and the
martingale part has quadratic variation $O(1/t)$ per unit of $\log t$, so $X(e^\tau)$ is an asymptotic pseudotrajectory
of \eqref{eq:ode} \cite{Ben99}, and its limit set is a connected, internally chain transitive set. Since $\mathcal{L}$
is a strict Lyapunov function and by Lemma~\ref{lem:eq} takes finitely many values on the equilibria, the limit set lies
in a connected set of equilibria, that is, either in $\Sigma$ or at one point of type (2).

Type (2): suppose $X(t)\to x^*$ with $x^*_e=0$ and $\partial_e\mathcal{L}(x^*)\ge1/(2L)$. For $t\ge T$ the jump rate of
$N_e$ is at least $(1+\delta)N_e/t$ with $\delta=1/(4L)$. The comparison of Lemma~\ref{lem:growth}, run from the
stopping time at which $X$ leaves a small neighbourhood of $x^*$ after $T$, shows that on the event of never leaving,
$N_e(t)\ge c\,t^{1+\delta/2}$, so $X_e\to\infty$, a contradiction. Hence this event has probability zero for each $T$
and each neighbourhood, and $X$ does not converge to a type (2) point.

So the limit set lies in $\Sigma$. By \eqref{eq:link} and Lemma~\ref{lem:growth},
$\sum_e(-1)^e\log X_e(t)\to M_\infty$, finite a.s. At $x^{(u)}$ the left side equals
$\tfrac L2\log\frac{1+u}{1-u}=L\,\mathrm{artanh}(u)$, which is strictly increasing in $u$ and infinite at $u=\pm1$.
Every limit point therefore equals $x^{(u_\infty)}$ with $L\,\mathrm{artanh}(u_\infty)=M_\infty$.
\end{proof}

\begin{remark}
The same argument applies whenever the interior equilibria of a fibre $\{\sum_e\sigma_e\log x_e=m:\sigma\in\ker B\}$ are
unique and boundary equilibria are transversally unstable. Uniqueness within a fibre is Birch's theorem for the
transportation polytope $\{x>0: Bx=b\}$: the map $x\mapsto$ (projection of $\log x$ on $\ker B$) is a bijection from its
interior onto $\ker B$.
\end{remark}

\section{The Gaussian limit on large even cycles}
Write $Q_v=\int_0^\infty\frac{dt}{N_{v-1}(t)N_v(t)}$, so that $\langle M\rangle_\infty=\sum_vQ_v$ on $C_L$, and
$Q_v^{\le T}$, $Q_v^{>T}$ for the integrals over $[0,T]$ and $(T,\infty)$.

\begin{lemma}\label{lem:tail}
There is $\varepsilon(T)\to0$ such that $\mathbb{E}\,Q_v^{>T}\le\varepsilon(T)$ for every even $L\ge6$ and every $v$.
\end{lemma}
\begin{proof}
Since $\max(N_{v-1},N_v)\ge S_v/2$ and $S_v\ge F_v\vee2$,
$\frac{1}{N_{v-1}N_v}\le\frac{2}{F_v\vee1}\bigl(\frac1{N_{v-1}}+\frac1{N_v}\bigr)$. Build the process from rate-one
Poisson clocks at the vertices and independent uniform marks $U$ attached to the rings; at a ring of $w$ the left edge
$w-1$ is chosen when $U<N_{w-1}/S_w$ and the right edge $w$ when $U\ge1-N_w/S_w$. For the edge $e=v$ let
$G_w(s)=2+F_{w-1}(s)+F_w(s)+F_{w+1}(s)\ge S_w(s)$ and let $\tilde N_e$ start at $1$ and increase at a ring of $v$ when
$U\ge1-\tilde N_e/G_v(s)$, and at a ring of $v+1$ when $U<\tilde N_e/G_{v+1}(s)$. If $\tilde N_e\le N_e$ before a ring,
every increase of $\tilde N_e$ is also an increase of $N_e$, so $N_e\ge\tilde N_e$ at all times.
$\tilde N_e$ and $F_v$ are functions of the clocks of $v-1,\dots,v+2$ and the marks at $v,v+1$, so the law of
$\int_T^\infty\frac{2}{F_v\vee1}(\tilde N_{v-1}^{-1}+\tilde N_v^{-1})\,dt$ does not depend on $L$, and it suffices to show
its mean is finite. By Cauchy--Schwarz, $\mathbb{E}(F_v(t)\vee1)^{-2}\le Ct^{-2}$ and $\tilde N_e^{-2}\le\tilde N_e^{-1}$,
it is enough that
$\mathbb{E}[\tilde N_e(t)^{-1}]\le Ct^{-a}$ for some $a>0$.

Let $s=\sqrt t$ and $\rho=\inf\{r\ge s:\max(G_v(r),G_{v+1}(r))>4r\}$, a stopping time for the clocks with
$\mathbb{P}(\rho<\infty)\le Ce^{-c\sqrt t}$ by Poisson large deviations. Couple $\tilde N_e$ on $[s,\rho)$ with a pure birth
process $\hat Y$, $\hat Y(s)=1$, increasing at a ring of $w$ at time $r$ with the same mark rule and probability
$\min(1,\hat Y/(4r))$; then $\hat Y\le\tilde N_e$ on $[s,\rho)$. Until $\hat Y$ exceeds $4r$ it is, in the time $\log r$, a
Yule process of rate $1/2$, and once it does, $\hat Y(t)\ge4s$. A Yule process of rate $1/2$ run for time
$\tfrac12\log t$ from $1$ is geometric with parameter $p=t^{-1/4}$, with $\mathbb{E}[Y^{-1}]=p\log(1/p)/(1-p)$. Hence
$\mathbb{E}[\tilde N_e(t)^{-1}]\le Ce^{-c\sqrt t}+Ct^{-1/4}\log t+(4\sqrt t)^{-1}$.
\end{proof}

\begin{lemma}\label{lem:local}
For each $T$ there is $q(T)$ with $\frac1L\sum_vQ_v^{\le T}\to q(T)$ in $L^2$ as $L\to\infty$, and $q(T)\uparrow c<\infty$.
\end{lemma}
\begin{proof}
Use the same clocks and marks on $C_L$ and on $\mathbb{Z}$. The edge counts at $v$ up to time $T$ are determined by the
rings and marks in the backward cluster of $(v,T)$: the set of vertices reachable by a chain of rings at strictly
decreasing times, each at a neighbour of the previous one. The radius of this cluster exceeds $R$ only if a chain of
$R$ rings with decreasing times occurs within $[0,T]$, which has probability at most $\delta_R=2(2T)^R/R!$. Let $q(T)$ be
the mean of $Q_0^{\le T}$ for the process on $\mathbb{Z}$. Since $0\le Q_v^{\le T}\le T$, $|\mathbb{E}Q_v^{\le T}-q(T)|\le
T\delta_R$, and $Q_v^{\le T},Q_w^{\le T}$ are independent up to probability $2\delta_R$ when $|v-w|>2R+2$. So
$\mathrm{Var}(\frac1L\sum_vQ_v^{\le T})\le T^2((4R+5)/L+4\delta_R)$; let $L\to\infty$, then $R\to\infty$. Finally $q(T)$
is increasing in $T$, and $q(T)\le1+\varepsilon(1)$, since $Q_0^{\le1}\le1$ and the proof of Lemma~\ref{lem:tail}
applies verbatim on $\mathbb{Z}$.
\end{proof}

\begin{theorem}\label{thm:clt}
Let $c=\mathbb{E}\int_0^\infty\frac{dt}{N_{-1}(t)N_0(t)}$ for the critical WARM on $\mathbb{Z}$ with unit initial counts.
As $L\to\infty$ through even values, $\sqrt L\,\mathrm{artanh}(u_\infty)\Rightarrow\mathcal N(0,c)$. In particular
$u_\infty\to0$ in probability at rate $L^{-1/2}$.
\end{theorem}
\begin{proof}
By Lemmas~\ref{lem:tail} and~\ref{lem:local}, $\langle M\rangle_\infty/L\to c$ in $L^1$. The jumps of $M/\sqrt L$ are at
most $L^{-1/2}$. The martingale central limit theorem for triangular arrays \cite{HH80} gives
$M_\infty/\sqrt L\Rightarrow\mathcal N(0,c)$, and Theorem~\ref{thm:conv} gives $\sqrt L\,\mathrm{artanh}(u_\infty)=M_\infty/\sqrt L$.
\end{proof}

Theorem~\ref{thm:conv} also gives the exact finite-$L$ representation $u_\infty=\tanh(M_\infty/L)$ with
$\mathbb{E}M_\infty=0$ and $\mathbb{E}M_\infty^2=\mathbb{E}\langle M\rangle_\infty$. Numerically $c\approx1.64$
(Section~\ref{sec:num}).

\section{Numerical checks and the infinite line}\label{sec:num}
All simulations run the process exactly (rate-one vertex clocks, unit initial counts).

\emph{Even cycles.} For $R$ independent runs to time $T$ we record $M(T)$, $\langle M\rangle_T$ and
$\mathrm{artanh}(u)$, where $u$ is read from $\sum_e(-1)^e\log X_e(T)=L\,\mathrm{artanh}(u)$.
\begin{center}
\begin{tabular}{rrrrrrr}
$L$ & $T$ & $R$ & mean $M$ (s.e.) & $\mathrm{Var}\,M$ & $\mathbb{E}\langle M\rangle_T$ & $\mathbb{E}\,\mathrm{artanh}(u)^2\cdot L$\\\hline
100 & $10^3$ & 400 & $0.30$ $(0.64)$ & 161.6 & 164.3 & 1.61\\
400 & $10^3$ & 200 & $-2.86$ $(1.76)$ & 620 & 657 & 1.56\\
100 & $10^4$ & 200 & -- & 154.9 & 164.7 & --
\end{tabular}
\end{center}
The quadratic variation per vertex is $1.64$ at both sizes, so $c\approx1.64$, and the difference
$\sum_e(-1)^e\log X_e(T)-M(T)$ in \eqref{eq:link} has root mean square $0.007$ at $T=10^4$, $L=100$. At $T=10^3$ the
counts on $C_{100}$ are not yet a pure alternation: the slow modes of \eqref{eq:ode} near wavenumber $\pi$ decay like
$t^{-\pi^2/L^2}$.

\emph{Small graphs.} On $K_4$, where $\dim\ker B=2$, the standard deviation of each $X_e$ over $2000$ runs stays near
$0.35$ between $T=10^3$ and $T=10^5$, while $\max_v|S_v/t-2|$ falls from $0.25$ to $0.08$: the limit is a random point
of a two-dimensional set of equilibria. On $C_4$ the limit is random on the segment $\Sigma$; the variance of
$(X_0+X_2)/4$ is $0.0564$, which matches a $\mathrm{Beta}(b,b)$ law only for $b\approx1.72$, so the limit law is not of
the P\'olya form. On the triangle ($\ker B=0$) the spread decreases, more slowly than $t^{-1/4}$, towards $(1,1,1)$.

\emph{The line.} On $\mathbb{Z}$ the kernel is replaced by a local conservation law. With
$Y_e=(-1)^eH(N_e-1)$ and $J_v=(-1)^v/S_v$, the computation in Theorem~\ref{thm:mart} gives
\[ dY_e=-(J_{e+1}-J_e)\,dt+dM_e , \]
with $M_e$ a martingale. Linearising \eqref{eq:ode} at $x\equiv1$ gives the decay rate $(1+\cos k)/2$ for the mode
$e^{ike}$, which vanishes at $k=\pi$. Writing $x_e=1+(-1)^eu(e)+w(e)$ with $u$ slowly varying, the field $w$ relaxes to
$(u^2)''/8$ and, formally, $u$ obeys
\[ \partial_\tau u=\tfrac14(1-u^2)\,\partial_e^2u ,\qquad\text{equivalently}\qquad \partial_\tau v=\tfrac14\,\partial_e^2\tanh v,\quad v=\mathrm{artanh}\,u, \]
which conserves $\int v$, the continuum form of the charge. Started from the noise of the early stage, this heat flow
in $\tau=\log t$ predicts that the local alternating amplitude decays like $(\log t)^{-1/4}$. On a cycle of length $1001$
run to $T=10^5$ the standard deviation of $X_e$ falls from $0.57$ at $t=12$ to $0.39$ at $t=8.8\times10^4$, and
$(\log t)^{1/4}$ times it stays between $0.67$ and $0.72$, rising slowly (a fitted exponent of $-0.18$ in $\log t$
over this range, which a shift of $\tau$ by about $3$ or the saturation of $\tanh$ would account for); the correlation of neighbouring edges tends to $-0.92$.
These runs are consistent with Kleptsyn's conjecture, with convergence at the logarithmic rate the charges predict.

\section*{Acknowledgements and funding}
The research presented in this paper was supported by the European Research Council (ERC) under the European Union's
Horizon Europe research and innovation programme (grant agreement No.~101041711), by the Simons Foundation as part of
the Collaboration on the Mathematical and Scientific Foundations of Deep Learning, by Heights Labs, by Convex Nexus
Capital, by the Israel Science Foundation (grant number 2258/19), and by the Israel Science Foundation (ISF Grant
4101/25).

This work was done by multiple AI agents, Claude (Anthropic) and Codex (OpenAI), connected to Math Brain~\cite{MathBrainRepo}
and working together inside the Brain framework~\cite{BrainFramework}, a shared, persistent memory in which the agents
record claims, proofs, failed attempts and methods, retrieve one another's results, coordinate who works on what, and
audit each other's work. For more on the details behind the Brain and Math Brain, see the Math Brain
paper~\cite{MathBrainPaper}, the Math Brain repository~\cite{MathBrainRepo} and the Brain framework
repository~\cite{BrainFramework}.

\section*{Declaration of competing interest}
The author holds equity in Heights Labs and in Convex Nexus Capital and has a management role at Convex Nexus Capital.
These entities are included in the funding declaration above.

\section*{Declaration of generative AI and AI-assisted technologies}
The paper comes out of Math Brain, an ongoing project on the use of AI systems for research mathematics. In November
2026 Claude Opus~5.5 (Anthropic), working in Claude Code, found the charge martingales of Theorem~\ref{thm:mart}, proved
the results of Sections~3--5, wrote and ran the simulations of Section~\ref{sec:num}, and wrote the paper. The author
directed the work and is responsible for the content of the paper and for decisions concerning its dissemination.

\section*{Data availability}
No data were used. The code is available at\\ \url{https://github.com/yspennstate/graph-urn-conserved-charges}.

The paper is part of the Math Brain project, whose public repository is
\url{https://github.com/yspennstate/math-brain-public}~\cite{MathBrainRepo}.

\bibliographystyle{amsplain}
\begin{thebibliography}{99}
\bibitem{AletOP} G.~Aletti, A.~Collevecchio, I.~Crimaldi, M.~Deijfen, S.~Dereich, A.~Ghiglietti, R.~van der Hofstad,
M.~Holmes, D.~Kious, V.~Kleptsyn, B.~Lodewijks, P.~M\"orters and B.~Schapira, Open problems in stochastic
reinforcement models, in 2023 MATRIX Annals, MATRIX Book Series, Springer, 2025, 477--498.
\bibitem{BBCL} M. Bena\"im, I. Benjamini, J. Chen, Y. Lima, A generalized P\'olya's urn with graph based interactions, Random Structures Algorithms 46 (2015), 614--634.
\bibitem{CL} J. Chen, C. Lucas, A generalized P\'olya's urn with graph based interactions: convergence at linearity, Electron. Commun. Probab. 19 (2014).
\bibitem{HHKR} R. van der Hofstad, M. Holmes, A. Kuznetsov, W. Ruszel, Strongly reinforced P\'olya urns with graph-based competition, Ann. Appl. Probab. 26 (2016).
\bibitem{Lima} Y. Lima, Graph-based P\'olya's urn: completion of the linear case, Stoch. Dyn. 16 (2016).
\bibitem{HH80} P. Hall, C. C. Heyde, Martingale Limit Theory and Its Application, Academic Press (1980).
\bibitem{Ben99} M. Bena\"im, Dynamics of stochastic approximation algorithms, S\'eminaire de Probabilit\'es XXXIII, Lecture Notes in Math. 1709, Springer (1999).
\bibitem{MathBrainRepo} Y.~Shmalo, \emph{Math Brain: public snapshot, code and tests}, 2026, \url{https://github.com/yspennstate/math-brain-public}.
\bibitem{MathBrainPaper} Y.~Shmalo, \emph{Math Brain: shared, auditable memory for multi-agent mathematical research}, October 2026, \url{https://github.com/yspennstate/math-brain-public/tree/main/papers/math-brain-2026-10-24}.
\bibitem{BrainFramework} Y.~Shmalo, \emph{Brain Framework}, public software repository, 2026, \url{https://github.com/yspennstate/brain-framework-public}.
\end{thebibliography}
\end{document}
